\newcounter{english}
\documentclass{article}

% packages
\usepackage{amsmath, amsthm, thmtools, amsfonts, amssymb, luacode, catchfile, hyperref, ifthen, graphicx, bbm}
\ifcsname c@kobocompile\endcsname
	\usepackage[a5paper, total={1072pt, 1448pt}, margin=10pt, includeheadfoot]{geometry} % set page margins
\else
	\usepackage[a4paper, margin=50pt, includeheadfoot]{geometry}
\fi
\usepackage[shortlabels]{enumitem}
\usepackage[skip=3pt, indent=0pt]{parskip}

% language
\usepackage[bidi=basic, layout=tabular, provide=*]{babel}
\ifcsname c@english\endcsname
	\babelprovide[main, import]{english}
\else
	\babelprovide[main, import, alph=letters.gershayim]{hebrew}
	\setlist[enumerate,2]{label={\alph*.}}
	\babelprovide{rl}
\fi
\babelfont{rm}[Renderer=Harfbuzz]{Libertinus Serif}
\babelfont{sf}{Libertinus Sans}
\babelfont{tt}{Libertinus Mono}

% style
\ifcsname c@stoppagebreaks\endcsname
\else
	\AddToHook{cmd/section/before}{\clearpage}	% Add line break before section
\fi
\linespread{1.3}
\setcounter{secnumdepth}{0}		% Remove default number tags from sections, this won't do well with theorems
\AtBeginDocument{\setlength{\belowdisplayskip}{3pt}}
\AtBeginDocument{\setlength{\abovedisplayskip}{3pt}}
\graphicspath{ {../images/} }

% operators
\DeclareMathOperator\cis{cis}
\DeclareMathOperator\Sp{Span}
\DeclareMathOperator\tr{tr}
\DeclareMathOperator\im{Im}
\DeclareMathOperator\re{Re}
\DeclareMathOperator\diag{diag}
\DeclareMathOperator*\lowlim{\underline{lim}}
\DeclareMathOperator*\uplim{\overline{lim}}
\DeclareMathOperator\rng{rng}
\DeclareMathOperator\Sym{Sym}
\DeclareMathOperator\Arg{Arg}
\DeclareMathOperator\Log{Log}
\DeclareMathOperator\dom{dom}
\DeclareMathOperator\supp{Supp}
\DeclareMathOperator\var{Var}
\DeclareMathOperator\cov{Cov}
\DeclareMathOperator\End{End}
\DeclareMathOperator\id{id}
\DeclareMathOperator\cl{cl}
\DeclareMathOperator\gal{Gal}
\DeclareMathOperator\dist{dist}
\DeclareMathOperator\form{form}
\DeclareMathOperator\term{term}
\DeclareMathOperator\sent{sent}
\DeclareMathOperator\Th{Th}
\DeclareMathOperator\Hom{Hom}
\DeclareMathOperator\Aut{Aut}
\DeclareMathOperator\Char{char}
\DeclareMathOperator\vol{vol}
\DeclareMathOperator\GL{GL}
\DeclareMathOperator\SL{SL}
\DeclareMathOperator\lcm{lcm}

% commands
\newcommand{\cupdot}{\mathbin{\mathaccent\cdot\cup}}
\newcommand{\indicator}{\mathbbm{1}}
\newcommand{\AAA}[0]{\mathbb{A}}
\newcommand{\Aa}[0]{\mathcal{A}}
\newcommand{\BB}[0]{\mathbb{B}}
\newcommand{\Bb}[0]{\mathcal{B}}
\newcommand{\CC}[0]{\mathbb{C}}
\newcommand{\Cc}[0]{\mathcal{C}}
\newcommand{\DD}[0]{\mathbb{D}}
\newcommand{\Dd}[0]{\mathcal{D}}
\newcommand{\EE}[0]{\mathbb{E}}
\newcommand{\FF}[0]{\mathbb{F}}
\newcommand{\Ff}[0]{\mathcal{F}}
\newcommand{\Ii}[0]{\mathcal{I}}
\newcommand{\Gg}[0]{\mathcal{G}}
\newcommand{\HH}[0]{\mathbb{H}}
\newcommand{\Hh}[0]{\mathcal{H}}
\newcommand{\Kk}[0]{\mathcal{K}}
\newcommand{\Ll}[0]{\mathcal{L}}
\newcommand{\Mm}[0]{\mathcal{M}}
\newcommand{\NN}[0]{\mathbb{N}}
\newcommand{\Nn}[0]{\mathcal{N}}
\newcommand{\PP}[0]{\mathbb{P}}
\newcommand{\Pp}[0]{\mathcal{P}}
\newcommand{\QQ}[0]{\mathbb{Q}}
\newcommand{\RR}[0]{\mathbb{R}}
\newcommand{\Rr}[0]{\mathcal{R}}
\newcommand{\Ss}[0]{\mathcal{S}}
\newcommand{\TT}[0]{\mathbb{T}}
\newcommand{\Tt}[0]{\mathcal{T}}
\newcommand{\UU}[0]{\mathbb{U}}
\newcommand{\Uu}[0]{\mathcal{U}}
\newcommand{\Vv}[0]{\mathcal{V}}
\newcommand{\Ww}[0]{\mathcal{W}}
\newcommand{\Xx}[0]{\mathcal{X}}
\newcommand{\Yy}[0]{\mathcal{Y}}
\newcommand{\ZZ}[0]{\mathbb{Z}}
\newcommand{\Zz}[0]{\mathcal{Z}}
\newcommand{\acts}[0]{\circlearrowright}
\newcommand{\explain}[2] {
	\begin{flalign*}
		 && \text{#2} && \text{#1}
	\end{flalign*}
}
\newcommand{\maketitleprint}[1][green]{
	\begin{center}
		\includegraphics[width=6cm]{dragon_#1}
	\end{center}
}

% theorem commands
\newtheoremstyle{c_remark}
	{}	% Space above
	{}	% Space below
	{}% Body font
	{}	% Indent amount
	{\bfseries}	% Theorem head font
	{}	% Punctuation after theorem head
	{.5em}	% Space after theorem head
	{\thmname{#1}\thmnumber{ #2}\thmnote{ \normalfont{\text{(#3)}}}}	% head content
\newtheoremstyle{c_definition}
	{3pt}	% Space above
	{3pt}	% Space below
	{}% Body font
	{}	% Indent amount
	{\bfseries}	% Theorem head font
	{}	% Punctuation after theorem head
	{.5em}	% Space after theorem head
	{\thmname{#1}\thmnumber{ #2}\thmnote{ \normalfont{\text{(#3)}}}}	% head content
\newtheoremstyle{c_plain}
	{3pt}	% Space above
	{3pt}	% Space below
	{\itshape}% Body font
	{}	% Indent amount
	{\bfseries}	% Theorem head font
	{}	% Punctuation after theorem head
	{.5em}	% Space after theorem head
	{\thmname{#1}\thmnumber{ #2}\thmnote{ \text{(#3)}}}	% head content

\ifcsname c@english\endcsname
	\theoremstyle{plain}
	\newtheorem{theorem}{Theorem}[section]
	\newtheorem*{theorem*}{Theorem}
	\newtheorem{lemma}[theorem]{Lemma}
	\newtheorem{proposition}[theorem]{Proposition}
	\newtheorem*{proposition*}{Proposition}

	\theoremstyle{definition}
	\newtheorem{definition}[theorem]{Definition}
	\newtheorem*{definition*}{Definition}
	\newtheorem{example}{Example}[section]
	\newtheorem{exercise}{Exercise}[section]

	\theoremstyle{remark}
	\newtheorem*{remark}{Remark}
	\newtheorem*{solution}{Solution}
	\newtheorem{corollary}[theorem]{Corollary}
	\newtheorem{notation}[theorem]{Notation}
\else
	\theoremstyle{c_plain}
	\newtheorem{theorem}{משפט}[section]
	\newtheorem*{theorem*}{משפט}
	\newtheorem{lemma}[theorem]{למה}
	\newtheorem{proposition}[theorem]{טענה}
	\newtheorem*{proposition*}{טענה}

	\theoremstyle{c_definition}
	\newtheorem{definition}[theorem]{הגדרה}
	\newtheorem*{definition*}{הגדרה}
	\newtheorem{example}{דוגמה}[section]
	\newtheorem{exercise}{תרגיל}[section]

	\theoremstyle{c_remark}
	\newtheorem*{remark}{הערה}
	\newtheorem*{solution}{פתרון}
	\newtheorem{corollary}[theorem]{מסקנה}
	\newtheorem{notation}[theorem]{סימון}
\fi

% Questions related commands
\newcounter{question}
\setcounter{question}{1}
\newcounter{sub_question}
\setcounter{sub_question}{1}

\ifcsname c@english\endcsname
	\newcommand{\question}[1][0]{
		\ifthenelse{#1 = 0}{}{\setcounter{question}{#1}}
		\section{Question \arabic{question}}
		\addtocounter{question}{1}
		\setcounter{sub_question}{1}
	}

	\newcommand{\subquestion}[1][0]{
		\ifthenelse{#1 = 0}{}{\setcounter{sub_question}{#1}}
		\subsection{Part \alph{sub_question}}
		\addtocounter{sub_question}{1}
	}
\else
	\newcommand{\question}[1][0]{
		\ifthenelse{#1 = 0}{}{\setcounter{question}{#1}}
		\section{שאלה \arabic{question}}
		\addtocounter{question}{1}
		\setcounter{sub_question}{1}
	}

	\newcommand{\subquestion}[1][0]{
		\ifthenelse{#1 = 0}{}{\setcounter{sub_question}{#1}}
		\subsection{סעיף \localecounter{letters.gershayim}{sub_question}}
		\addtocounter{sub_question}{1}
	}
\fi

% import lua and start of document
\directlua{common = require ('../common')}

\GetEnv{AUTHOR}

% headers
\author{\AUTHOR}
\date\today

\title{Space of analytic functions over unit disk --- Notes}
\setcounter{secnumdepth}{2}
% chktex-file 9
% chktex-file 17

\usepackage{fancyhdr}
\pagestyle{fancy}
\renewcommand{\headrulewidth}{0pt}

\begin{document}
\maketitle
\maketitleprint[blue]{}

\tableofcontents

\section{Lecture 1 --- 18.10.2026}
The main object of this course is a Hardy Space, for $0 < p < \infty$, the Hardy space $\HH^p$ consists of all holomorphic functions $f : \DD \to \CC$ such that,
\[
	\sup_{0 < r < 1} \int_{- \pi}^{\pi} {| f(r e^{i t}) |}^p\ dt < \infty
\]
and,
\[
	\HH^\infty
	= \{ \text{bounded analytic functions on } \DD \}
\]
for $p < q$ it is clear that $\HH^p \supseteq \HH^q \supseteq \ldots \supseteq BMO \supseteq \HH^\infty$.
We study boundary behavior, in particular Fatou's theory; for a.e. $\theta \in [0, 2 \pi)$, a function $f \in \HH^\infty$ has a non-tangential limit at $e^{i \theta}$,
\[
	z_n to e^{i \theta}
	\implies \lim_{n \to \infty} f(z_n) \text{ exists}
\]
in particular $\lim_{r \to 1^-} f(r e^{i \theta})$ exists.

Another question we study is Zero sets, a set of points $\{ a_n \} \subseteq \DD$ is the Zero set of an $\HH^\infty$ function if and only if $\sum (1 - | a_n |) < \infty$,
in other words it converges to the boundary quickly (Blaschke condition).

We also study Invariant subspaces $z X \subseteq X$, if $z(z) \in X$ then $z \cdot f(z) \in X$ (Beurling, 1949).

Interpolating sequences. A sequence $\{ a_n \} \subseteq \DD$ is interpolating if for any $b_n \in l^\infty$ there exists $f \in \HH^\infty$ such that $f(a_n) = b_n$.
The answer was given by Carleson: a sequence is interpolating if and only if it is separated and has Carleson condition.

Finally, Corona problem: given $f_1, f_2 \in \HH^\infty$ such that $|f_1(z)| + |f_2(z)| > \delta$ for all $z \in \DD$, then there are $g_1, g_2 \in \HH^\infty$ such that $f_1 g_1 + f_2 g_2 = 1$.
The Space $BMO$ is going to be involved with the solution to this problem.
There are also Bergman spaces, $A^p = \{ f : \DD \xrightarrow{\text{hol}} \CC \mid \int_{\DD} {| f(z) |}^p\ dA < \infty \}$ and Dirichlet space,
\[
	\Dd
	= \{ f : \DD \xrightarrow{\text{hol}} \CC \mid \int_{\DD} {| f'(z) |}^2\ dA < \infty \}
\]
The textbook is ''$H_p$ spaces'' by P. Koosis. The preliminaries are Complex analysis 1, Measure theory, Hilbert spaces and little bit of Functional analysis.

Reminder about $\Ll^p$ spaces, ${(\Ll^p)}^* \simeq \Ll^q$ for $1 \le p < \infty$ where $\frac{1}{p} + \frac{1}{q} = 1$.
Then $\lVert g \rVert_{L^q} = \sup_{\lVert f \rVert_{L^p}} \int f \cdot g\ dt$. Another fact is that $C_c(\RR) \subseteq \Ll^p(\RR)$ is dense.
$(T_t f)(x) = f(x + t)$ is continuous in $\Ll^p$ in the sense that $t \to T_t f$ is continuous,
\[
	\lVert T_t f - f \rVert_{L^p} \xrightarrow{t \to 0} 0
\]
the proof is left as exercise.
$C {(\partial \DD)}^* \simeq \Mm(\partial \DD)$ and $C_0 {(\RR)}^* \simeq \Mm(\RR)$.
We can define the measure $\mu$ that takes $f \to \int_{\partial \DD} f\ d \mu$.

\subsection{Harmonic functions}
\begin{definition}[Harmonic function]
	a function $u : \Omega \to \RR$ is harmonic if one of the following conditions holds:
	\begin{enumerate}
		\item $u$ is $C^2$ and $\Delta u = \frac{\partial^2}{\partial x^2} u + \frac{\partial^2}{\partial y^2} = 0$
		\item $u \in C$ satisfies the mean-value property, namely for any disk $B(z_0, r) \subseteq \Omega$ it holds that,
			\[
				u(z_0)
				= \frac{1}{2 \pi} \int_{- \pi}^{\pi} u(z_0 + r e^{i \theta})\ d \theta
			\]
		\item $u$ is locally the real part of an analytic function.
	\end{enumerate}
\end{definition}
\begin{example}
	$\log |z|$ is harmonic in $\CC \setminus \{ 0 \}$.
\end{example}
We would have wanted to take $f = \log z$ but we would've encountered a problem with continuity.
Note that condition 3 holds for the imaginary part as well.
For the sake of convenience we will treat condition 3 as the main definition in this course.
Suppose that $u : B(0, R) \to \RR$ is harmonic, then $u = \re f$ for some analytic function $f$, but any such function has an analytic series,
\[
	f(z)
	= \sum_{n = 0}^\infty a_n z^n
\]
such that,
\[
	\liminf_{n \to \infty} \frac{1}{\sqrt[n]{|a_n|}}
	\ge R
\]
Taking real parts, we get,
\[
	u(r r^{i \theta})
	= \sum_{n = -\infty}^\infty A_n r^{|n|} e^{i n \theta}
\]
In order for this to be real we need $A_{-n} = \overline{A_n}$ and the condition,
\[
	\forall n \in \ZZ,\ 
	\liminf_{n \to \infty} \frac{1}{\sqrt[|n|]{|A_n|}}
	\ge R
\]

\subsection{Poisson integral formula}
Poisson formula gives us some generalization of the mean-value property in the sense that each point is valued by an harmonic function as the boundary of a consisting ball, up to weight function (kernel).
Let $u : B(0, R) \to \RR$ with $R > 1$ and $r < 1$ then,
\[
	u(r e^{i \theta})
	= \frac{1}{2 \pi} \int_{- \pi}^{\pi} u(e^{i t}) \left[ \sum_{- \infty}^\infty r^{|n|} e^{i n (\theta - t)} \right]\ dt
\]
We know that,
\[
	\frac{1}{2 \pi} \int_{- \pi}^{\pi} e^{i n t}\ dt
	= \begin{cases}
		0 & n \ne 0 \\
		1 = 0
	\end{cases}
\]
and that if $u$ is defined as a sequence, then,
\[
	u(r e^{i t})
	= \frac{1}{2 \pi} \int_{- \pi}^{\pi} \sum_{- \infty}^\infty A_n e^{i n t} r^{|n|} e^{i n (\theta - t)}\ dt
\]
thus we get the exact formula.
\begin{definition}[Poisson kernel]
	For any $\theta$ we define,
	\[
		P_r(\theta)
		= \sum_{- \infty}^\infty r^{|n|} e^{i n \theta}
		= \frac{1 - r^2}{1 + r^2 - 2 r \cos \theta}
		= \re(\frac{1 + r e^{i \theta}}{1 - r e^{i \theta}})
	\]
\end{definition}
Then,
\[
	u(r e^{i \theta})
	= \frac{1}{2 \pi} \int_{- \pi}^{\pi} u(e^{i t}) P_r(\theta - t)\ dt
\]
\begin{proposition}
	The following properties hold for Poisson kernel:
	\begin{enumerate}
		\item $P_r(\theta) > 0$
		\item $P_r$ is $2 \pi$-periodic
		\item $\frac{1}{2 \pi} \int_{- \pi}^\pi P_r(\theta)\ d \theta = 1$
		\item For any $\delta > 0$, $P_r(\theta) \to 0$ uniformly for $\pi > | \theta | > \delta$, as $r \to 1$
		\item $P_r(\theta)$ is even
	\end{enumerate}
\end{proposition}
\begin{remark}
	We can also denote using Poisson kernel with,
	\[
		u(e r^{i \theta})
		= \frac{1}{2 \pi} u(e^{i t}) P_r(\theta - t)\ dt
	\]
\end{remark}

Suppose that $f \in \Ll^p(\partial \DD)$ or equivalently $\in \Ll^p([-\pi, \pi])$, and let $u = P[f]$ Poisson extension of $f$, namely,
\[
	u(r e^{i t})
	= \frac{1}{2 \pi} \int_{- \pi}^{\pi} f(t) P_r(\theta - t)\ dt
\]
then $u$ is harmonic in $\DD$.
We can also consider $u + i \tilde{u}$ for $\tilde{u}$ the harmonic conjugate.
So we can plug $\Ll^p$ functions into the PIF (Poisson Integration Formula) to obtain harmonic functions in the disk.
What kind of functions do we get by doing so, and how are $u$ and $f$ related? Is $f \to u \to \tilde{u} \to \tilde{f}$ bounded in $\Ll^p$ (Singular integral operators)?
That is a delicate question which we will answer during the course.
\begin{theorem}[Extension to harmonic functions]
	Suppose that $u(r e^{i \theta})$ is a harmonic function on the unit disk. If $1 < p < \infty$ and,
	\[
		\sup_{0 < r < 1} \int_{|z| = r} {| u(r e^{i \theta}) |}\ d \theta < \infty
	\]
	then there is a function $f : [-\pi, \pi] \to \RR$ in $L^p$ such that $u(r e^{i \theta}) = \frac{1}{2 \pi} \int_{- \pi}^{\pi} P_r(\theta - t) f(t)\ dt$.
\end{theorem}
\begin{proof}
	Compactness:
	We know that ${(\Ll^q)}^* \simeq \Ll^p$. Choose a countable dense set $G_1, G_2, \ldots \in \Ll^q$.
	Define,
	\[
		\Ll G_j
		= \lim_{j \to \infty} \frac{1}{2 \pi} \int_{- \pi}^{\pi} G_k(\theta) u((1 - \frac{1}{n_j}) e^{i \theta})\ d \theta
	\]
	$G_k \in \Ll^q$ and $u((1 - n_j) e^{i \theta})$ is bounded because $L^p$ norm is bounded, then the whole integral is bounded.
	Remember that $|\int f g| \le \lVert f \rVert_p \cdot \lVert g \rVert_q$ for any $f \in L^p, g \in L^q$.
	For a general $G \in \Ll^q$, we approximate $G$ by $G_{k_l} \in \Ll^q$ and define $\Ll G = \lim_{l \to \infty} \Ll G_{k_l}$.
	It follows that there exists a function $f \in \Ll^p$ such that $\Ll G = \int_{-\pi}^{\pi} f G$.
	Fix $-\pi \le \theta \le \pi$ and choose $G(t) = P_r(\theta - t)$, let $v = P[f]$,
	\[
		\Ll P_r(\theta - t)
		= \int_{-\pi}^{\pi} f(t) P_r(\theta - t)\ dt
		= 2 \pi v(r e^{i \theta})
		= \lim_{j \to \infty} \int_{-\pi}^{\pi} P_r(\theta - t) u((1 - \frac{1}{n_j}) e^{i \theta})
	\]
	The dilate $u((1 - \frac{1}{n_j}) e^{i \theta})$ is harmonic on $B(0, \frac{1}{1 - \frac{1}{n_j}}) \subseteq B(0, 1)$ therefore we can use Poisson formula,
	\[
		\lim_{j \to \infty} \int_{-\pi}^{\pi} P_r(\theta - t) u((1 - \frac{1}{n_j}) e^{i \theta})
		= \lim_{j \to \infty} 2 \pi u((1 - \frac{1}{n_j}) r e^{i \theta})
		= 2 \pi \cdot u(r e^{i \theta})
	\]
\end{proof}
\begin{theorem}[Existence of measure to harmonic functions]
	Suppose that $u(e r^{i \theta})$ is a harmonic function on $\DD$ such that,
	\[
		\sup_{0 < r < 1} \int_{-\pi}^{\pi} |u(r e^{i \theta})|\ d \theta < \infty
	\]
	then there exists a measure $\mu$ such that,
	\[
		u(r e^{i \theta})
		= \frac{1}{2 \pi} \int_{-\pi}^{\pi} P_r(\theta - t)\ d \mu
	\]
\end{theorem}
\begin{example}
	Remember that $P_r(\theta) = \frac{1 - r^2}{1 + r^2 - 2 \cos \theta}$, in the case ``$r = 1$'' acts like $\delta_1$.
\end{example}

\section{Lecture 2 --- 25.10.2026}
\subsection{Measures and total variation}
A \textbf{signed measure} is the difference of two positive finite measures, $\mu = \mu_+ - \mu_-$.
The \textbf{total variation measure} $|\mu| = \mu_+ + \mu_-$, this leads to the identity,
\[
	\left\lvert \int f\ d \mu \right\rvert
	\le \int |f|\ d |\mu|
\]
We define $\lVert \mu \rVert_M = |\mu|(X)$.
A complex measure is, in similar manner, a sum of two measures, $\mu = \lambda + i \nu$ where $\lambda, \nu$ are signed measures.
In this case we define the total variation as $|\mu| = |\lambda| + |\nu| = \lambda_- + \lambda_- + \nu_+ + \nu_-$.
\begin{theorem}[Banach-Alaoglu]
	Let $X$ be a Banach space, The unit ball in $X^*$ is weak-$*$ compact.
\end{theorem}
\begin{example}
	${(\Ll^q)}^* \simeq \Ll^p$ for $1 < \le \infty$ and $\frac{1}{p} + \frac{1}{q} = 1$.
\end{example}
Suppose $f_1, f_2, \ldots$ with $\sup \lVert f_j \rVert_{L^p} < \infty$.
$\{ f_j \}$ may not have a continuous subsequence in $\Ll^p$, BA (Banach-Alaoglu) implies that there are $j_k$ such that $f_{j_k} \to f$ in weak-$*$,
\[
	\forall g \in \Ll^q,
	\quad
	\int f_{j_k} g\ d x
	\to \int f g\ d x
\]

\subsection{Harmonic functions in R-disk}
Denote $\DD_R = \{ z \in \CC \mid |z| < R \}$.
Function $u : \DD_R \to \RR$ is harmonic if,
\[
	u(r e^{i \theta})
	= \sum_{n = -\infty}^\infty A_n r^{|n|} e^{i n \theta}
\]
where $A_{-n} = \bar{A}_n$ and radius $\ge R$.

We can consider complex-valued harmonic functions, in this case we can omit the first condition and say that $u$ is analytic iff $A_{-n} = 0$ for $n > 0$,
\[
	u(r e^{i \theta})
	= \sum_{n \ge 0} A_n r^n e^{i n \theta}
\]
$u$ is \textbf{anti-analytic} iff $A_n = 0$ for $n > 0$.
Any harmonic function $h$ can be written as $h = i + i v$ where $u, v$ are real-valued harmonic functions, but also as $h = f + \bar{g}$ where $f, g$ are analytics. \\
We say that $v$ is a \textbf{harmonic-conjugate} of $u$ if $u + i v$ is analytic.
\begin{lemma}
	$v$ is unique up to an additive constant.
\end{lemma}
If $v_1, v_2$ are such conjugates then $(u + i v_1) - (u + i v_2) = i(v_1 - v_2)$ which is constant. \\
Existence:
\[
	\tilde{u}(r e^{i \theta})
	= \sum_{n = 0\infty}^\infty (-i \sgn n) A_n r^{|n|} e^{i n \theta}
\]
From this definition it follows that $\tilde{A}_n = -i A_n$ for $n > 0$ and so $\tilde{A}_{-n} = i A_{-n}$, meaning that indeed $\tilde{A}_{-n} = \bar{\tilde{A}}_n$.

When discussing Poisson integral formula we assumed that $R > 1$,
\[
	u(r e^{i \theta})
	= \frac{1}{2 \pi} \int_{-\pi}^{\pi} u(e^{i t}) \left[ \sum_{n = -\infty}^\infty r^{|n|} e^{i n (\theta - t)} \right]\ dt
\]
When $R > 1$, $|A_n| \le C \Theta^{-n}$ for some $\Theta > 1$.
\begin{example}
	The series $z + z^2 + \cdots$ has a radius $1 = \limsup \sqrt[|n|]{|A_n|}$, to enlarge it we have to set the coefficients to be much smaller.
	$\frac{z}{2} + \frac{z^2}{4} + \frac{z^3}{8}$ has a radius of $2$.
\end{example}
We will define Poisson kernel as,
\[
	P_r(\theta)
	= \sum_{n = -\infty}^\infty r^{|n|} e^{i n \theta}
	= \frac{1 - r^2}{1 + r^2 - 2 r \cos \theta}
	= \re \frac{1 + r e^{i \theta}}{1 - r e^{i \theta}}
\]
Poisson kernel $P(r e^{i \theta})$ is a harmonic function in $\DD$.
This is clear as it can be written as a special case of the series defining a harmonic function, or as the real part of an analytic function.
In general we can achieve the harmonic conjugate with the formula,
\[
	\tilde{u}(r e^{i \theta})
	= \frac{1}{2 \pi} \int_{-\pi}^{\pi} u(e^{i t}) \left[ \sum_{n = -\infty}^\infty -i \sgn(n) r^{|n|} e^{i n (\theta - t)} \right]\ dt
\]
So we define the conjugate Poisson kernel as,
\[
	Q_r(\theta)
	= \sum_{n = -\infty}^\infty -i \sgn(n) r^{|n|} e^{i n \theta}
	= \frac{2 r \sin \theta}{1 + r^2 - 2 r \cos \theta}
	= \im \frac{1 + r e^{i \theta}}{1 - r r^{i \theta}}
\]

\subsection{Harmonic functions with bounded Lp means}
\begin{theorem}[Representation with Lp fuctions]
	Suppose $u$ is harmonic function on the unit disk.
	If,
	\[
		\sup_{0 < r < 1} \int_{-\pi}^{\pi} {|u(r e^{i \theta})|}^p\ d \theta < \infty
		\tag{1}
	\]
	then,
	\[
		u(r e^{i \theta})
		= \frac{1}{2 \pi} \int_{-\pi}^{\pi} F(t) P_r(\theta - t)\ dt
	\]
	where $F \in \Ll^p$.
\end{theorem}
\begin{proof}
	Set $F_n(\theta) = u((1 - \frac{1}{n}) e^{i \theta})$ for $\theta \in [-\pi, \pi]$.
	$(1)$ sways that $\lVert F_n \rVert_{L^p} \le C$ for all $n$.
	Banach-Alaoglu means that for any $G \in L^q$,
	\[
		\int_{-\pi}^{\pi} F_{n_k}(\theta) G(\theta)\ d \theta
		\to \int_{-\pi}^{\pi} F(\theta) G(\theta)\ d \theta
	\]
	Fix $0 \le r < 1$ and $t \in [-\pi, \pi]$ and take $G(\theta) = P_r(t - \theta)$, this is a smooth function.
	\[
		2 \pi u(r e^{i t})
		= 2 \pi \lim_{k \to \infty} u((1 - \frac{1}{n_k}) r e^{i t})
		= \lim_{k \to \infty} \int_{-\pi}^{\pi} u((1 - \frac{1}{n_k}) e^{i \theta}) P_r(t - \theta)\ d \theta
		= \int_{-\pi}^{\pi} F(\theta) P_r(t - \theta)\ d \theta
	\]
\end{proof}
\begin{theorem}[equivalency in the case of 1]
	Suppose $u$ is harmonic function on the unit disk and $p = 1$.
	If,
	\[
		\sup_{0 < r < 1} \int_{-\pi}^{\pi} {|u(r e^{i \theta})|}\ d \theta < \infty
	\]
	then,
	\[
		u(r e^{i \theta})
		= \frac{1}{2 \pi} \int_{-\pi}^{\pi} P_r(\theta - t)\ d \mu(t)
	\]
	where $\mu$ is a signed measure.
\end{theorem}
\begin{corollary}
	Suppose $u : \DD \to \RR$ is a positive harmonic function.
	Then $u = P[\mu]$ for $\mu$ a positive measure, namely $d u = P d \mu$.
\end{corollary}
\begin{example}
	$P(r e^{i \theta}) = \frac{1 - r^2}{1 + r^2 - 2 r \cos \theta}$ is a positive harmonic function,
	\[
		P(r e^{i \theta})
		= \frac{1}{2 \pi} \int_{-\pi}^{\pi} P_r(\theta - t) d (2 \pi \delta_0(t))
	\]
	It follows that $P = P[2 \pi \delta_0]$.
\end{example}
\begin{example}[Plateau's problem]
	Suppose that $\gamma \subseteq \RR^3$ is a closed smooth path,
	find a surface $S \subseteq \RR^3$ with $\partial S = \gamma$ of minimal area.

	Let $m = \inf\{ \vol_2 S \mid \partial S = \gamma \}$, for any $n \in \NN$ there is a surface $S_n$ with $\partial S_n = \gamma$ and $\vol_2 S_n < m + \frac{1}{n}$.
\end{example}
\begin{theorem}
	Suppose $F(\theta) \in \Ll^p$, where $1 < p \le \infty$, and $u = P[F]$. Then,
	\[
		\sup_{0 < r < 1} \int_{-\pi}^{\pi} {| u(r e^{i \theta})|}^p\ d \theta < \infty
	\]
	\begin{proof}
		Set $F_r(\theta) = u(r e^{i \theta})$ for $\theta \in [-\pi, \pi]$.
		Then there is a function $G \in \Ll^q$ with $\lVert G \rVert_{L^q} = 1$ such that,
		\begin{align*}
			\lVert F_r \rVert_{L^p}
			& = \int_{-\pi}^{\pi} F_r(\theta) G(\theta)\ d \theta \\
			& = \left\lvert \frac{1}{2 \pi} \int_{-\pi}^{\pi} \int_{-\pi}^{\pi} F(\theta - t) P_r(t)\ dt\ G(\theta)\ d \theta \right\rvert \\
			& \le \frac{1}{2 \pi} \int_{-\pi}^{\pi} \lVert F \rVert_{L^p} \lVert G \rVert_{L^q} P_r(t)\ dt \\
			& = \lVert F \rVert_{L^p} \lVert G \rVert_{L^q} \frac{1}{2 \pi} \int_{-\pi}^{\pi} P_r(t)\ dt \\
			& = \lVert F \rVert_{L^p} \lVert G \rVert_{L^q} \\
			& = \lVert F \rVert_{L^p}
		\end{align*}
		In the case $p = 1$, $G \in \Ll^\infty$ and $G = \frac{\bar{F}_r}{|F_r|}$, the inequalities are a bit different,
		\begin{align*}
			\lVert F_r \rVert_{L^1}
			& = \int_{-\pi}^{\pi} F_r(\theta) G(\theta)\ d \theta \\
			& = \frac{1}{2 \pi} \int_{-\pi}^{\pi} \int_{-\pi}^{\pi} P_r(\theta - t) d \mu(t) G(\theta)\ d \theta \\
			& \le \frac{1}{2 \pi} \int_{-\pi}^{\pi} \int_{-\pi}^{\pi} P_r(\theta - t) |d \mu(t)|\ d \theta \\
			& = \int_{-\pi}^{\pi} |d \mu(\theta)| \\
			& = \lVert \mu \rVert_M
		\end{align*}
		and the claim holds.
	\end{proof}
\end{theorem}

\subsection{Elementary study of boundary behavior}
\begin{theorem}[Extension of unit function to boundary]
	Suppose $F$ is a continuous function on $[-\pi, \pi]$ and $u = P[F]$.
	Then $u$ extends continuously to $\overline{\DD}$ and,
	\[
		\lim_{z \to e^{i \theta}} u(z)
		= F(\theta)
	\]
\end{theorem}
Notice that $P_r(\theta) > 0$, is even, $2 \pi$-periodic, $\int_{-\pi}^{\pi} P_r(\theta)\ d \theta = 2 \pi$
and $\forall \delta > 0 \exists r_0,\ r > r_0 \implies \int_{|t| > \delta} P_r(\theta)\ d \theta < \delta$.
\begin{proof}
	We show $\lim_{r \to 1} u(r e^{i \theta}) = F(\theta)$.
	\[
		u(r e^{i \theta})
		= \frac{1}{2 \pi} \int_{-\pi}^{\pi} F(\theta - t) P_r(t)\ d t
	\]
	Then,
	\[
		u(r e^{i \theta}) - F(\theta)
		= \frac{1}{2 \pi} \int_{-\pi}^{\pi} \int_{-\pi}^{\pi} F(\theta - t) P_r(t)\ dt
	\]
	and in particular,
	\[
		|u(r e^{i \theta}) - F(\theta)|
		\le \frac{1}{2 \pi} \int_{-\pi}^{\pi} \int_{-\pi}^{\pi} |F(\theta - t) - F(\theta)| P_r(t)\ dt
		= \frac{1}{2 \pi} (\int_{|t| < \delta} + \int_{|t| > \delta})
	\]
	We show that these two integrals are bounded by $\frac{\varepsilon}{2}$.
	Since $F$ is uniformly continuous, there exists a $\delta > 0$ such that,
	\[
		|F(\theta - t) - F(\theta)| < \frac{\varepsilon}{2}
	\]
	for $|t| < \delta$.
	Denote the two integrals by $I$ and $I\!I$, then,
	\[
		|I|
		\le \frac{1}{2 \pi} \int_{|t| < \delta} \frac{\varepsilon}{2} P_r(t)\ dt
		\le \frac{\varepsilon}{2}
	\]
	Since $F$ is continuous, it is bounded, $M = \lVert F \rVert_{L^\infty} < \infty$.
	\[
		|I\!I|
		\le \frac{1}{2 \pi} \int_{|t| > \delta} 2 M P_r(t)\ dt
	\]
	but by Poisson kernel properties, we can choose a value such that $r$ is close to $1$.
\end{proof}
\begin{theorem}
	Suppose $F(\theta) \in \Ll^p$ with $1 < p < \infty$ and $u = P[F]$.
	Set $F_r(\theta) = u(r e^{i \theta})$, then $F_r \to F$ in $L^p$ as $r \to 1$.
\end{theorem}
\begin{proof}
	\[
		\lVert F_r(\theta) - F(\theta) \rVert_{L^p}
		= \left\lVert \frac{1}{2 \pi} \int_{-\pi}^{\pi} (F(\theta - t) - F(\theta)) P_r(t) \right\rVert_{L^p}
		\le \frac{1}{2 \pi} \int_{-\pi}^{\pi} \lVert F(\theta - t) - F(\theta) \rVert_{L^p} P_r(t)\ dt
		\le \frac{1}{2 \pi} (\int_{|t| < \delta} + \int_{|t| > \delta})
	\]
	and we once again denote them $I$ and $I\!I$.
	$I < \frac{\varepsilon}{2}$ by the same claim as the last theorem.
	If $\lVert F \rVert_{L^p} = M$ then $\lVert F(\theta - t) - F(\theta) \rVert_{L^p} \le 2 M$ hence $I\!I < \frac{\varepsilon}{2}$ by a choice of $\delta$.
\end{proof}

\listoftheorems[title=theorems \& definitions,ignoreall,show={theorem,definition},swapnumber,onlynamed={proposition}]
\end{document}
